\item Experimenta e interpreta algunas manifestaciones y aplicaciones de la electricidad, e identifica los cuidados que requiere su uso al revisar los protocolos de seguridad.
\item Relaciona e interpreta fenómenos comunes del magnetismo y experimenta con la interacción entre imanes.
\item Experimenta e interpreta el comportamiento de la luz como resultado de la interacción entre electricidad y magnetismo.
\item Explica el funcionamiento de aparatos tecnológicos de comunicación, a partir de las ondas electromagnéticas.
\item Indaga algunos avances recientes en la comprensión sobre la evolución del Universo y su composición.
\item Indaga cómo se lleva a cabo la exploración de los cuerpos celestes, por medio de la detección y procesamiento de las ondas electromagnéticas que emiten.
\item Relaciona e interpreta las características y dinámica del Sistema Solar con la gravitación y el movimiento de los planetas, en particular el caso de la Tierra y la Luna.
\item Formula hipótesis que relacionan la actividad humana con el aumento de temperatura en el planeta y la emisión de gases de efecto invernadero; diferencia entre calor, radiación y temperatura al explicar los procesos que lo originan.
\item Indaga sobre fenómenos meteorológicos extremos como olas de calor, ciclones tropicales, sequías y lluvias torrenciales; representa y explica su distribución en el mundo.
\item Propone medidas de mitigación y adaptación, encaminadas al cuidado del medio ambiente y el bienestar común, viables para su aplicación en su escuela y comunidad.